\section{Original geometry and positive cohomology across a scheme isomorphism}

Fix an actual scheme isomorphism $e:X\simeq Y$.  Every pullback below
uses its actual forward map, so a module on $Y$ pulls back to a module
on $X$.

\begin{lemma}[The actual module unit and quasicoherence across an isomorphism]
\label{lem:kn-scheme-iso-module-unit}
\leanok
\lean{EnvelopeProperty.Vanishing.SchemeIsoGeometry.pullbackPushforwardUnitIso,EnvelopeProperty.Vanishing.SchemeIsoGeometry.pullback_isQuasicoherent}
For every original module $M$ on $Y$, the actual adjunction unit gives
$M\simeq e_*e^*M$.  If $M$ is quasicoherent, so is the same actual
$e^*M$.
\end{lemma}
\begin{proof}
Pullback along a scheme isomorphism is an equivalence.  Its actual
right adjoint is pushforward, so the unit of this adjunction is an
isomorphism at every original module.  An isomorphism is an open
immersion, and its pullback comparison with open restriction identifies
$e^*M$ with the original restriction of $M$.  Restriction preserves
quasicoherence, proving the second assertion.
\end{proof}

\begin{lemma}[Whole positive original cohomology vanishing is invariant]
\label{lem:kn-scheme-iso-positive-cohomology}
\leanok
\lean{EnvelopeProperty.Vanishing.SchemeIsoGeometry.pullback_subsingleton_H_iff}
\uses{lem:kn-scheme-iso-module-unit}
Let $p:Y\to\operatorname{Spec}R$ have an actual projective
factorization and let $M$ be quasicoherent.  For every $q>0$,
\[
 H^q(X,e^*M)=0\quad\Longleftrightarrow\quad H^q(Y,M)=0.
\]
\end{lemma}
\begin{proof}
The projective factorization gives separatedness of $Y$ and an actual
finite ordered affine cover.  The original inverse images under $e$
form an affine cover of $X$.  Both original quasicoherent modules are
acyclic on the actual affine intersections.  The native ordered affine
Cech pushforward comparison therefore identifies vanishing of whole
positive cohomology of $e^*M$ with that of its actual pushforward.
The preceding whole module unit identifies this pushforward with the
same original $M$.  Applying the actual cohomology functor to this
unit gives the stated equivalence of vanishing.
\end{proof}

\begin{lemma}[The original geometric ample datum pulls back]
\label{lem:kn-scheme-iso-ample-pullback}
\leanok
\lean{EnvelopeProperty.Vanishing.SchemeIsoGeometry.isAmple_pullback,EnvelopeProperty.Vanishing.SchemeIsoGeometry.isAmple_pullback_of_triangle}
\uses{lem:fc-original-ample-closed-restriction}
If the original line $L$ is geometrically ample over $p$, the actual
line $e^*L$ is geometrically ample over $p\circ e$.
If $p\circ e=p_X$, this is ampleness over the same specified $p_X$.
\end{lemma}
\begin{proof}
Take the original positive power $r$ and its actual closed projective
embedding exhibiting $L^{\otimes r}$ as the pullback of the original
$\mathcal O(1)$.  The composite embedding with $e$ is closed because
$e$ is an isomorphism.  Native pullback composition and the actual
comparison $e^*(L^{\otimes r})\simeq(e^*L)^{\otimes r}$ give precisely
the same embedding criterion for $e^*L$, with the original positive
$r$ unchanged.  Its invertibility is preserved by actual pullback.
Transporting the structural map along the given triangle gives the
last assertion.
\end{proof}

\begin{lemma}[Original exterior and canonical modules across a compatible isomorphism]
\label{lem:kn-scheme-iso-exterior-pullback}
\leanok
\lean{EnvelopeProperty.Vanishing.SchemeIsoGeometry.inv_structure,EnvelopeProperty.Vanishing.SchemeIsoGeometry.exteriorPullbackIso,EnvelopeProperty.Vanishing.SchemeIsoGeometry.canonicalPullbackIso}
\uses{sc0-canonical-port:can:relative-exterior-cotangent-sheaf-open-restriction}
Let $f_X:X\to\operatorname{Spec}k$ and $f_Y:Y\to\operatorname{Spec}k$
be the specified original structural maps over a field, with
$f_Y\circ e=f_X$.  For every $i$ there is an actual module isomorphism
\[
 e^*\Omega^i_{Y/k}\simeq\Omega^i_{X/k}.
\]
In degree $d$ this is the comparison of the same original canonical
objects.  The inverse isomorphism has the inverse structural triangle
$f_X\circ e^{-1}=f_Y$.
\end{lemma}
\begin{proof}
Use the actual restriction-versus-pullback comparison for the open
immersion $e$, followed by the native open-immersion exterior
cotangent comparison.  Its target is the actual exterior module
relative to the composite $f_Y\circ e$.  Transport that structural map
along the given triangle, obtaining the original $f_X$-relative
module.  The canonical object is by definition this same actual
exterior construction in the specified degree $d$.  Composing the
original structural triangle with $e^{-1}$ proves the inverse one.
\end{proof}

\begin{corollary}[The same original powers and canonical twists are retained]
\label{cor:kn-scheme-iso-hodge-tensor-pullback}
\leanok
\lean{EnvelopeProperty.Vanishing.SchemeIsoGeometry.exteriorTensorPullbackIso,EnvelopeProperty.Vanishing.SchemeIsoGeometry.hodgeTensorPowerPullbackIso,EnvelopeProperty.Vanishing.SchemeIsoGeometry.canonicalTwistPullbackIso,EnvelopeProperty.Vanishing.SchemeIsoGeometry.canonicalTensorObjPullbackIso}
\uses{lem:kn-scheme-iso-exterior-pullback}
For every original module $L$ and every $m,i$,
\[
 e^*(L^{\otimes m}\otimes\Omega^i_{Y/k})
 \simeq(e^*L)^{\otimes m}\otimes\Omega^i_{X/k}.
\]
The unpowered exterior twist and the original canonical-first twist
also have their whole actual pullback comparisons; in particular,
$e^*(\omega_Y\otimes L)\simeq\omega_X\otimes e^*L$ when $d$ is the
specified smooth relative dimension.
\end{corollary}
\begin{proof}
The actual strong monoidal pullback comparison separates the two
original tensor factors.  Apply the native power-pullback isomorphism
to the first factor and the preceding actual exterior comparison to
the second.  The unpowered and canonical-first statements use the
same tensor comparison in the displayed original order of factors.
The native comparison between the monoidal tensor and the explicitly
sheafified coherent tensor gives the same whole comparison for that
original coherent tensor as well.
\end{proof}
